problem:  
Let \(M^k \subset \mathbb{H}^{n+1}\) be a complete, properly embedded minimal submanifold with smooth, closed asymptotic boundary \(\Lambda = \partial_\infty M \subset S^n_\infty\).  
In the Fefferman–Graham compactification \((\bar{\mathbb{H}}^{n+1}, g_{\mathbb{H}^{n+1}} = \rho^{-2}(d\rho^2 + h_{S^n}))\), \(M\) can be described near infinity by a normal graph \(u(\theta,\rho)\) over the asymptotic cone \(C_\Lambda\)\!,
\[
u(\theta,\rho) = \rho\, u_1(\theta) + \rho^2 u_2(\theta) + \cdots,
\]
so that the induced metric on \(M\) takes the asymptotic form
\[
g_M = \rho^{-2}\big(h_0 + \rho^2 h_2 + \cdots + \rho^k h_k + \mathcal{O}(\rho^{k+1})\big),
\]
where \(h_0\) is the induced metric on \(\Lambda\).

The coefficients \(h_{2j}\) in this *ambient extension* are recursively determined by the intrinsic and extrinsic curvatures of \(\Lambda\), except that at order \(k\) a formally undetermined, trace-free tensor \(h^{(k)}\) appears—the analogue of the holographic “obstruction tensor.”

**Conjecture (Asymptotic Obstruction Energy):**  
For each integer \(k\ge2\), define a *boundary energy* functional on smooth, closed submanifolds \(\Lambda \subset S^n\) by
\[
\mathcal{E}_k(\Lambda) := \int_{\Lambda} \!\! \langle h^{(k)}_{\text{TF}},\, h^{(k)}_{\text{TF}} \rangle_{h_0} \, d\mu_{h_0},
\]
where \(h^{(k)}_{\text{TF}}\) is the trace-free part of the formally undetermined coefficient in the expansion above.

Then:

1. **(Existence)** \(\mathcal{E}_k(\Lambda)\) is finite for any smooth closed \(\Lambda\subset S^n\) bounding a properly embedded minimal \(M\subset \mathbb{H}^{n+1}\);  
2. **(Independence)** if \(M_1\) and \(M_2\) are two minimal submanifolds with the same asymptotic boundary \(\Lambda\), then \(\mathcal{E}_k(\Lambda)\) computed from either agrees;  
3. **(Characterization)** \(\mathcal{E}_k(\Lambda) = 0\) if and only if \(\Lambda\) is a minimal submanifold of \(S^n_\infty\);  
4. **(Conformal invariance)** under conformal transformations of \(S^n_\infty\), \(\mathcal{E}_k(\Lambda)\) transforms with a universal scaling, defining a conformally invariant quadratic form on the moduli of boundary embeddings.

Equivalently, does the tensor \(h^{(k)}_{\text{TF}}\) give a canonical geometric measure of the asymptotic obstruction to total geodesicity of \(M\), yielding a natural *renormalized curvature energy* on the boundary \(\Lambda\)?

---

Why it is a “good” problem:  

1. **Structurally rooted:**  
   Unlike ad hoc energy definitions, this conjecture draws directly from the *known formal expansion structure* of minimal submanifolds in hyperbolic space—precisely where an undetermined asymptotic term \(h^{(k)}\) first appears, reflecting the intrinsic–extrinsic obstruction in the Fefferman–Graham recursion.  
   Extending the scalar “obstruction tensor” phenomenon from ambient metric theory to submanifold geometry aligns perfectly with the established analytic machinery, making the problem intrinsically grounded.

2. **Analytic and geometric bridge:**  
   The conjecture seeks to connect PDE asymptotics of the minimal surface equation (analysis of the obstruction term) with *conformal geometric invariants* on the boundary, paralleling how the ambient obstruction tensor connects Poincaré–Einstein metrics to conformal curvature quantities.  
   It thereby unites geometric analysis, conformal geometry, and holographic renormalization in a precise mathematical statement.

3. **Clarity and attainability:**  
   Each ingredient—formal expansions, the leading undetermined tensor, trace-free decomposition—already has a well-defined meaning. The open question is whether this data yields a *geometrically natural, filling-independent energy*. This is a concrete, plausible research conjecture with testable predictions (vanishing ⇔ boundary minimality).

4. **Simplicity with profound implications:**  
   The statement is conceptually simple: the first undetermined term in the expansion of a hyperbolic minimal submanifold should encode a canonical conformal energy on the boundary. Its resolution would illuminate the analytic structure controlling boundary regularity and potentially define new conformally invariant geometrical energies beyond the Willmore functional.

5. **Potential for field integration:**  
   Solving or even partially understanding this would likely require harmonic analysis, variational methods, and representation theory of curvature tensors within ambient conformal geometry—precisely the kind of cross-disciplinary integration that typifies important advances in modern geometric analysis.

In summary, this problem is mathematically coherent, logically consistent, and conceptually aligned with existing structures while posing a genuinely open and deep conjectural link between two major domains—*minimal submanifold asymptotics and conformal boundary geometry*.  
It thus satisfies the criteria of a “good,” research-level mathematical problem.